\documentclass[10pt, letterpaper]{article}

% Packages:
\usepackage[
    ignoreheadfoot, % set margins without considering header and footer
    top=1.3 cm, % seperation between body and page edge from the top
    bottom=1.3 cm, % seperation between body and page edge from the bottom
    left=1.6 cm, % seperation between body and page edge from the left
    right=1.6 cm, % seperation between body and page edge from the right
    footskip=0.6 cm, % seperation between body and footer
    % showframe % for debugging 
]{geometry} % for adjusting page geometry
\usepackage{titlesec} % for customizing section titles
\usepackage{tabularx} % for making tables with fixed width columns
\usepackage{array} % tabularx requires this
\usepackage[dvipsnames]{xcolor} % for coloring text
\definecolor{primaryColor}{RGB}{0, 79, 144} % define primary color
\usepackage{enumitem} % for customizing lists
\usepackage{fontawesome5} % for using icons
\usepackage{amsmath} % for math
\usepackage[
    pdftitle={Razeen Wasif's CV},
    pdfauthor={Razeen Wasif},
    pdfcreator={LaTeX with RenderCV},
    colorlinks=true,
    urlcolor=primaryColor
]{hyperref} % for links, metadata and bookmarks
\usepackage[pscoord]{eso-pic} % for floating text on the page
\usepackage{calc} % for calculating lengths
\usepackage{bookmark} % for bookmarks
\usepackage{lastpage} % for getting the total number of pages
\usepackage{changepage} % for one column entries (adjustwidth environment)
\usepackage{paracol} % for two and three column entries
\usepackage{ifthen} % for conditional statements
\usepackage{needspace} % for avoiding page brake right after the section title
\usepackage{iftex} % check if engine is pdflatex, xetex or luatex

% Ensure that generate pdf is machine readable/ATS parsable:
\ifPDFTeX
    \input{glyphtounicode}
    \pdfgentounicode=1
    % \usepackage[T1]{fontenc} % this breaks sb2nov
    \usepackage[utf8]{inputenc}
    \usepackage{lmodern}
\fi



% Some settings:
\AtBeginEnvironment{adjustwidth}{\partopsep0pt} % remove space before adjustwidth environment
\pagestyle{empty} % no header or footer
\setcounter{secnumdepth}{0} % no section numbering
\setlength{\parindent}{0pt} % no indentation
\setlength{\topskip}{0pt} % no top skip
\setlength{\columnsep}{0cm} % set column seperation
\makeatletter
\let\ps@customFooterStyle\ps@plain % Copy the plain style to customFooterStyle
\patchcmd{\ps@customFooterStyle}{\thepage}{
    \color{gray}\textit{\small Razeen Wasif - Page \thepage{} of \pageref*{LastPage}}
}{}{} % replace number by desired string
\makeatother
\pagestyle{customFooterStyle}

\titleformat{\section}{\needspace{4\baselineskip}\bfseries\large}{}{0pt}{}[\vspace{1pt}\titlerule]

\titlespacing{\section}{
    % left space:
    -1pt
}{
    % top space:
    0.2 cm
}{
    % bottom space:
    0.1 cm
} % section title spacing

\renewcommand\labelitemi{$\circ$} % custom bullet points
\newenvironment{highlights}{
    \begin{itemize}[
        topsep=0.03 cm,
        parsep=0.03 cm,
        partopsep=0pt,
        itemsep=0pt,
        leftmargin=0.4 cm + 10pt
    ]
}{
    \end{itemize}
} % new environment for highlights

\newenvironment{highlightsforbulletentries}{
    \begin{itemize}[
        topsep=0.03 cm,
        parsep=0.03 cm,
        partopsep=0pt,
        itemsep=0pt,
        leftmargin=10pt
    ]
}{
    \end{itemize}
} % new environment for highlights for bullet entries


\newenvironment{onecolentry}{
    \begin{adjustwidth}{
        0.2 cm + 0.00001 cm
    }{
        0.2 cm + 0.00001 cm
    }
}{
    \end{adjustwidth}
} % new environment for one column entries

\newenvironment{twocolentry}[2][]{
    \onecolentry
    \def\secondColumn{#2}
}{
    \hfill \secondColumn
    \endonecolentry
} % new environment for two column entries

\newenvironment{header}{
    \setlength{\topsep}{0pt}\par\kern\topsep\centering\linespread{1.25}
}{
    \par\kern\topsep
} % new environment for the header

\newcommand{\placelastupdatedtext}{% \placetextbox{<horizontal pos>}{<vertical pos>}{<stuff>}
  \AddToShipoutPictureFG*{% Add <stuff> to current page foreground
    \put(
        \LenToUnit{\paperwidth-2 cm-0.2 cm+0.05cm},
        \LenToUnit{\paperheight-1.0 cm}
    ){\vtop{{\null}\makebox[0pt][c]{
        \small\color{gray}\textit{}\hspace{\widthof{}}
    }}}%
  }%
}%

% save the original href command in a new command:
\let\hrefWithoutArrow\href

% new command for external links:
\renewcommand{\href}[2]{\hrefWithoutArrow{#1}{\ifthenelse{\equal{#2}{}}{ }{#2 }\raisebox{.15ex}{\footnotesize \faExternalLink*}}}

% =========================================================
% =========================================================
% =========================================================
% =========================================================
% =========================================================
% =========================================================

\begin{document}
    \newcommand{\AND}{\unskip
        \cleaders\copy\ANDbox\hskip\wd\ANDbox
        \ignorespaces
    }
    \newsavebox\ANDbox
    \sbox\ANDbox{}

    \placelastupdatedtext
    \begin{header}
        \textbf{\fontsize{24 pt}{24 pt}\selectfont Razeen Wasif}

        \vspace{0.1 cm}

        \normalsize
        
        \mbox{{\color{black}\footnotesize\faMapMarker*}\hspace*{0.13cm}Canberra}%
        \kern 0.25 cm%
        \AND%
        \kern 0.25 cm%
        \mbox{\hrefWithoutArrow{mailto:razeen.wasif66@gmail.com}{\color{black}{\footnotesize\faEnvelope[regular]}\hspace*{0.13cm}razeen.wasif66@gmail.com}}%
        \kern 0.25 cm%
        \AND%
        \kern 0.25 cm%
        \mbox{\hrefWithoutArrow{tel:0450004709}{\color{black}{\footnotesize\faPhone*}\hspace*{0.13cm}0450004709}}%
        \kern 0.25 cm%
        \AND%
        \kern 0.25 cm%
        \mbox{\hrefWithoutArrow{https://www.linkedin.com/in/razeenwasif}{\color{black}{\footnotesize\faLinkedinIn}\hspace*{0.13cm}Razeen Wasif}}%
        \kern 0.25 cm%
        \AND%
        \kern 0.25 cm%
        \mbox{\hrefWithoutArrow{https://github.com/razeenwasif}{\color{black}
        {\footnotesize\faGithub}\hspace*
        {0.13cm}razeenwasif}}%
        \kern 0.25 cm%
        \AND%
        \kern 0.25 cm%
        \mbox{\hrefWithoutArrow{https://razeenwasif.web.app}{\color{black}{\footnotesize\faLink}\hspace*
        {0.13cm}razeenwasif.web.app}}%
        

    \end{header}

    \vspace{0.1 cm}
    \begin{center}
        \textbf{\large Summary}
    \end{center}

    \begin{onecolentry}
        \begin{center}
            Software Engineer \& ML Researcher with a strong foundation in \textbf{AI, machine learning,} and \textbf{software construction}. My background in physics provides a rigorous analytical and quantitative approach to problem-solving. Experienced in designing and implementing AI models, building \textbf{training infrastructure}, and shipping \textbf{full-stack platforms} using \textbf{Python, C++, CUDA,} and \textbf{TypeScript}.
        \end{center}
    \end{onecolentry}

    \columnratio{0.64,0.36}
    \begin{paracol}{2}
    \setlength{\columnsep}{0.6 cm}

    \section{Professional Experience}

    \begin{twocolentry}{
        Feb 2026 -- Present
    }
        \textbf{ANU Virtuous Loop} \\
        \textit{Full-Stack Developer \& UX Lead}
    \end{twocolentry}

    \vspace{0.10 cm}
    \begin{onecolentry}
        \begin{highlights}
            \item Collaborating with an industry client to architect a scalable feedback dashboard on the \textbf{Microsoft Power Platform (Power Apps, Power BI, Dataverse)}.
            \item Spearheading a complete \textbf{UI/UX overhaul} (user research + wireframing), lifting navigation efficiency and data accessibility by 40\%.
            \item Engineering automated \textbf{Power Automate} pipelines to replace legacy manual processes with real-time feedback sync.
            \item Implementing \textbf{Row-Level Security (RLS)} and enterprise data validation to meet ANU's privacy and compliance standards.
        \end{highlights}
    \end{onecolentry}

    \section{Projects}

    \begin{twocolentry}{
        \href{https://prism-automl.web.app/}{Live} / \href{https://github.com/razeenwasif/Prism}{Github}
    }
        \textbf{Prism: GPU-Accelerated AutoML \& Entity Resolution Engine}
    \end{twocolentry}
    
    \vspace{0.05 cm}
    \begin{onecolentry}
        \begin{highlights}
            \item Engineered a \textbf{GPU-first AutoML pipeline} for heterogeneous tabular data — automated EDA, target-aware cleaning, and \textbf{Optuna}-powered hyperparameter optimisation.
            \item Integrated a \textbf{RAPIDS}-accelerated record linkage workflow for entity resolution, using async execution and semantic embeddings to maximise candidate pair recall.
            \item Developed custom CUDA-backed XGBoost and PyTorch wrappers with automated CUDA/MPS/CPU fallback, served via a \textbf{FastAPI + Docker + React/Vite} stack and an Ollama LLM for NL data exploration.
        \end{highlights}
    \end{onecolentry}

    \begin{twocolentry}{
        \href{https://github.com/razeenwasif/flux}{Github}
    }
        \textbf{Flux: AI-Native Desktop Browser}
    \end{twocolentry}

    \vspace{0.05 cm}
    \begin{onecolentry}
        \begin{highlights}
            \item Built an \textbf{AI-native desktop browser} in \textbf{Rust + Tauri v2} (native OS webviews, no Electron) — an \textbf{8--15 MB} binary with $<$80 MB idle RAM, using a \textbf{zero-copy DOM IPC pipeline} that shares page snapshots as \texttt{Arc<[u8]>} across the terminal, agent, and embedder.
            \item Integrated a fully \textbf{local LLM agent} (Gemma 12B via Ollama / llama.cpp) that reads the page DOM and performs preview-and-approve actions behind a destructive-action deny-list — no cloud, no token cost.
            \item Shipped a real-\textbf{PTY} terminal sidebar with a DOM-aware CLI, webview hibernation under memory pressure, a Brave-engine ad/tracker blocker, and an \textbf{AES-256-GCM} password vault with OS-keychain key storage.
        \end{highlights}
    \end{onecolentry}

    \begin{twocolentry}{
        \href{https://github.com/razeenwasif/omni}{Github}
    }
        \textbf{Omni: From-Scratch Hybrid Search Engine}
    \end{twocolentry}

    \vspace{0.05 cm}
    \begin{onecolentry}
        \begin{highlights}
            \item Hand-rolled a search engine in \textbf{std-only Rust} (Go crawler, Python eval harness) — inverted index, \textbf{BM25F} ranking, Block-Max WAND pruning, Porter stemming, and a hand-written \textbf{HNSW} ANN index, with no Lucene / Tantivy.
            \item Fused lexical and \textbf{768-dim semantic} retrieval via Reciprocal Rank Fusion with passage-level scoring, reaching \textbf{0.666 nDCG@10} (vs.~0.477 lexical-only) and \textbf{0.97 success@10} on a 12k-document graded benchmark.
            \item Engineered live incremental indexing with content-addressed segments, a background tiered-merge thread, and atomic \texttt{Arc} snapshot swaps (zero downtime); serves a custom \textbf{HTTP/1.0 + SSE} API with local-LLM RAG answers that powers the Flux omnibox.
        \end{highlights}
    \end{onecolentry}

    \begin{twocolentry}{
        \href{https://nexus-cluster.web.app}{Live}
    }
        \textbf{Nexus: Distributed Training Console}
    \end{twocolentry}

    \vspace{0.05 cm}
    \begin{onecolentry}
        \begin{highlights}
            \item Architected a standalone control plane and real-time monitoring console for distributed ML training, with its own backend, brand, and deployment surface.
            \item Built a \textbf{Firestore-backed} pod inventory and run-lifecycle model, surfacing NCCL liveness probes and per-rank telemetry rolled up live in a monochrome glass shell.
            \item Designed a streaming event console for multi-rank orchestration — per-step throughput, gradient norms, and rank-level failures without polling — on \textbf{React + TypeScript + Firebase} with a Python ingestion service.
        \end{highlights}
    \end{onecolentry}

    \begin{twocolentry}{
        \href{https://oracle-neuro-sym.web.app/}{Live}
    }
        \textbf{Oracle: Multi-Domain Identification Platform}
    \end{twocolentry}

    \vspace{0.05 cm}
    \begin{onecolentry}
        \begin{highlights}
            \item Engineered a live multi-domain identification platform that productizes the \textbf{BioCLIP 2.5} PlantCLEF model — image upload, top-k predictions with taxonomy context, served via a \textbf{React + TypeScript + Firebase} front end.
            \item Implemented the inference path with \textbf{4$\times$4 tile decomposition} and multi-resolution (224 / 336 px) ensembling, with logit-adjusted adaptive thresholding tuned for long-tail multi-label outputs.
            \item Validated end-to-end on the \textbf{PlantCLEF 2026} benchmark (7,806-class long-tail), reaching \textbf{0.418 Macro F1} on the held-out vegetation-plot test set.
        \end{highlights}
    \end{onecolentry}

    \begin{twocolentry}{
        \href{https://github.com/razeenwasif/onyx}{Onyx} / \href{https://github.com/razeenwasif/audiopulse}{AudioPulse} / \href{https://github.com/razeenwasif/boxtube-tui}{BoxTube}
    }
        \textbf{Terminal Workbench: Onyx · AudioPulse · BoxTube}
    \end{twocolentry}

    \vspace{0.05 cm}
    \begin{onecolentry}
        \begin{highlights}
            \item Designed and shipped three day-in-the-life TUIs in a four-day sprint — markdown notes, music, and video — each rendering its native media inside the terminal.
            \item \textbf{Onyx} (\textbf{Rust} / Ratatui): Obsidian-style markdown vault with \texttt{[[wikilinks]]} + backlinks, command palette, full-vault content search, ASCII graph view, and vim-style modal editing.
            \item \textbf{AudioPulse} (\textbf{Go} / Bubble Tea): Spotify-style music player playing full songs through an embedded \textbf{librespot} device (PKCE OAuth via Spotify Web API), album art rendered as 24-bit half-blocks.
            \item \textbf{BoxTube} (\textbf{Python} / Textual): YouTube client with full video playback rendered inside the terminal via \textbf{mpv} and the \textbf{kitty graphics protocol} (sixel / unicode fallbacks), no Google Data API required.
        \end{highlights}
    \end{onecolentry}

    \begin{twocolentry}{
        \href{https://github.com/razeenwasif/Comp-Architecture/tree/main/16-bit-cpu}{github}
    }
        \textbf{Custom 16-Bit RISC CPU Design \& Implementation}
    \end{twocolentry}

    \vspace{0.05 cm}
    \begin{onecolentry}
        \begin{highlights}
            \item Designed and simulated a \textbf{16-bit RISC processor} from basic logic gates — custom ISA, multi-stage datapath (ALU, register file, microprogrammed control unit), and a 16-bit address bus with dedicated PC / IR.
            \item Validated end-to-end through unit tests of individual components and full assembly-level programs (Fibonacci, array traversal).
        \end{highlights}
    \end{onecolentry}

    \switchcolumn


    \section{Education}

        
        \begin{onecolentry}
            \textbf{Australian National University}

            \textit{Bachelor of Advanced Computing (Honors)}
        \end{onecolentry}

        \vspace{0.05 cm}
        \begin{onecolentry}
            \begin{highlights}
                \item \textbf{Coursework:} Machine Learning, Artificial Intelligence, Computer Vision, Software Engineering, Computer Architecture, Algorithms and Data structures, Data analysis
            \end{highlights}
        \end{onecolentry}

        \vspace{0.05 cm}

    % =========================================================
    % =========================================================
    % =========================================================


    \section{Open Source Contributions}

        \begin{onecolentry}
            \textbf{Pygame Library} (\hrefWithoutArrow{https://github.com/pygame/pygame}{GitHub})
        \end{onecolentry}
        \vspace{0.05 cm}
        \begin{onecolentry}
            \begin{highlights}
                \item Identified and resolved a \textbf{segmentation fault} in the \texttt{PixelArray} module — debugged memory management in the \textbf{C backend} to prevent invalid access during array slicing.
            \end{highlights}
        \end{onecolentry}

    % =========================================================
    % =========================================================
    % =========================================================

   \section{Research}

        \begin{twocolentry}{
        \href{https://github.com/razeenwasif/Council}{Github}
        }
            \textbf{Quantized Domain-Specialist Verifiers for Multi-Agent LLM Scientific Reasoning}
        \end{twocolentry}

        \vspace{0.05 cm}
        \begin{onecolentry}
            \begin{highlights}
                \item Investigating a hybrid architecture where frontier API models drive multi-agent debate (\textbf{Council} `/discover`) and locally-hosted, fine-tuned, aggressively-quantized open-weights models act as a \textbf{domain-specialist verification layer} for domain-specific claims.
                \item Built a \textbf{single-GPU LoRA + DoRA + QLoRA pipeline} (Unsloth, NEFTune, sequence packing, 8-bit paged AdamW) on Mistral-7B-Instruct-v0.3 — producing physics / math / CS specialists trained partly on \textbf{Council-generated synthetic debate data}.
                \item Designing a controlled ablation across six bit-depths (FP16 / Q8 / Q6 / Q4 / Q3 / Q2) comparing \textbf{PTQ vs.~QAT} in the specialist-as-verifier regime; target venue: NeurIPS ENLSP / ICLR ME-FoMo / ICML AI4Science 2026.
            \end{highlights}
        \end{onecolentry}

    \vspace{0.1 cm}

        \begin{twocolentry}{
        \href{https://arm-wision.github.io/}{Project} / \href{https://www.kaggle.com/competitions/plantclef2026}{Kaggle}
        }
            \textbf{Partial Fine-Tuning of BioCLIP 2.5 for Multi-Label Plant ID (PlantCLEF 2026)}
        \end{twocolentry}

        \vspace{0.05 cm}
        \begin{onecolentry}
            \begin{highlights}
                \item Partial fine-tune of \textbf{BioCLIP 2.5 ViT-H/14} with auxiliary \textbf{genus / family / order / class} heads, injecting taxonomic structure into the loss to stabilise training on a 7,806-class long-tail benchmark.
                \item Engineered a \textbf{4$\times$4 tile inference} pipeline ensembled across \textbf{224 and 336-pixel} resolutions, with \textbf{logit-adjusted adaptive thresholding} per class to handle multi-label vegetation plots.
                \item Achieved \textbf{0.418 Macro F1} on the PlantCLEF 2026 leaderboard; paper submitted to \textbf{LifeCLEF / CLEF 2026} (ANU at PlantCLEF 2026).
            \end{highlights}
        \end{onecolentry}

    \vspace{0.1 cm}

        \begin{twocolentry}{
        \href{https://github.com/razeenwasif/TinyML-Quantization_Induced_GW_SNR_Degradation}{github}
        }
            \textbf{Signal-Preserving TinyML for Gravitational Wave Triggers}
        \end{twocolentry}

        \vspace{0.05 cm}
        \begin{onecolentry}
            \begin{highlights}
                \item Investigated the impact of \textbf{INT8 quantization} on phase-coherence in Gravitational Wave (GW) signals for deployment on ultra-low-power \textbf{ARM Cortex-M} hardware.
                \item Developed a "Signal-Preserving" framework using \textbf{Quantization-Aware Training (QAT)} and Hardware-Aware Distillation to minimize SNR loss in stochastic astrophysical data.
                \item Engineered custom \textbf{C++ inference kernels} bypassing TensorFlow Lite Micro for 4096 Hz strain data within a \textbf{<128KB SRAM} budget, achieving over 99\% data compression at the target FAR.
            \end{highlights}
        \end{onecolentry}
        
    % =========================================================
    % =========================================================
    % =========================================================

    

    
    \section{Technologies}

        \begin{onecolentry}
            % Rust, Ruby, Haskell
            \textbf{Languages:} Python, C++, CUDA, Java, TypeScript, ARM ASM, SQL \end{onecolentry}

        \vspace{0.1 cm}

        \begin{onecolentry}
            \textbf{ML / AI:} PyTorch, TensorFlow, Optuna, RAPIDS, BioCLIP, DINOv3 \end{onecolentry}

        \vspace{0.1 cm}

        \begin{onecolentry}
            \textbf{Systems:} Linux Kernel, TensorRT, NVIDIA DALI, FastAPI, Docker \end{onecolentry}

        \vspace{0.1 cm}

        \begin{onecolentry}
            \textbf{Platforms:} Microsoft Power Platform, Power BI, Dataverse, Firebase, Figma, React \end{onecolentry}

    % =========================================================
    % =========================================================
    % =========================================================
    % =========================================================
    % =========================================================
    % =========================================================

    \section{Certifications}
        
        
            \begin{onecolentry}
                \textbf{Deloitte Australia Technology Job Simulation}\newline
                \textit{Forage} \hfill \textit{Aug 2024}
            \end{onecolentry}
            \begin{onecolentry}
                \vspace{0.05 cm}
                \begin{itemize}
                    \item Data analysis for Deloitte's Technology team — Tableau dashboard, client proposal, Excel-based classification and business conclusions.
                \end{itemize}
            \vspace{0.05 cm}
            \end{onecolentry}

    % =========================================================
    % =========================================================
    % =========================================================
    % =========================================================
    % =========================================================
    % =========================================================



    \end{paracol}

\end{document}